% Section content template for individual Educative lessons
% This file contains the actual content for a specific section
% Generated from Educative API section content

% Section: Code for the Online Blackjack Game
% Section ID: 5678759047069696
% Section Slug: code-for-the-online-blackjack-game
% Generated: 2025-12-12T12:55:26.591222

Using various UML diagrams, we've covered different aspects of the Blackjack game and observed the attributes attached to the problem. Let us now explore the more practical side of things, where we will work on implementing the Blackjack game using multiple languages. This is usually the last step in an object-oriented design interview process.
\\
We have chosen the following languages to write the skeleton code of the different classes present in the Blackjack game:

\begin{itemize}
\item
\protect\phantomsection\label{jr_IYwlFb5nLMm9d5MNWO}
Java
\item
\protect\phantomsection\label{UWn1ZTPLFyB05xlluPOuH}
C\#
\item
\protect\phantomsection\label{FPsF7TqNPs4EKX206N2cj}
Python
\item
\protect\phantomsection\label{hR0MUOnnytnUYXhePr1TF}
C++
\item
\protect\phantomsection\label{qbFFkvlc-M8q-Kd8U1Rdb}
JavaScript
\end{itemize}
\\
If you want to skip directly to the implementation and see the complete workflow, simply scroll down or click~\hyperref[Executable-Code-Online-Blackjack-Game]{Executable Code: Online Blackjack Game}.

\subsection{Blackjack game classes}\label{c1sKl9D6pmmSnzINYEn7T}
\\
This section will provide the skeleton code of the classes designed in the class diagram lesson.
\begin{quote}
\textbf{Note:} For simplicity, we are not defining getter and setter functions. The reader can assume that all class attributes are private, accessed through their respective public getter methods, and modified only through their public method functions.
\end{quote}
\subsubsection{Enumerations and custom data type}\label{GriG8ycdNvZ6AOmRI7YuW}
\\
The following code defines the enumeration and custom data type used in the Blackjack game.
\\
\textbf{\texttt{Suit}:} We need to create an enumeration to keep track of the card's suit, whether diamond, spade, heart, or club.
\\
\textbf{\texttt{AccountStatus}:} We need to create an enumeration to keep track of the account's status, whether it is active, canceled, closed, blocked, or none.
\\
The \texttt{Person} class is used as a custom data type. The implementation of the \texttt{Person} class can be found below:
\begin{quote}
\textbf{Note:} JavaScript does not support enumerations, so we will use the \texttt{Object.freeze()} method as an alternative that freezes an object and prevents further modifications.
\end{quote}

\textbf{Definition of enums and custom datatypes}

\begin{lstlisting}[language=Java]
public enum Suit {
    HEART, SPADE, CLUB, DIAMOND
}

public enum AccountStatus {
    ACTIVE, CLOSED, CANCELED, BLACKLISTED, NONE
}

public class Person {
    private String name;
    private String streetAddress;
    private String city;
    private String state;
    private String country;

    // Constructor(s)
    // Getters and setters
}
\end{lstlisting}

\subsubsection{Card}\label{mn3zfNeJYir-RBT9xIj9B}
\\
This class contains the playing cards or cards used in the Blackjack game.

\textbf{The Card class}

\begin{lstlisting}[language=Java]
public class Card {
    private Suit suit;
    private int faceValue;

    public Card(Suit suit, int faceValue) { /* ... */ }
    public int getValue() { /* ... */ }
    public Suit getSuit() { /* ... */ }
    public int getFaceValue() { /* ... */ }
    @Override
    public String toString() { /* ... */ }
}
\end{lstlisting}

\subsubsection{Deck and Shoe}\label{d6DcwCJ0jKUAiX5iUpOvw-2}
\\
\texttt{Shoe} is a device to hold multiple \texttt{Deck}, and a \texttt{Deck} has 52 cards of four suits. One suit contains nine number cards (2--10) and four face cards (King, Queen, Jack, and Ace).

\textbf{The Deck and Shoe classes}

\begin{lstlisting}[language=Java]
import java.util.*;

public class Deck {
    private List<Card> cards;

    public Deck() { /* ... */ }
    public List<Card> getCards() { /* ... */ }
}

public class Shoe {
    private List<Deck> decks;

    public Shoe(int numberOfDecks) { /* ... */ }
    public void shuffle() { /* ... */ }
    public Card dealCard() { /* ... */ }
    public List<Deck> getDecks() { /* ... */ }
}
\end{lstlisting}

\subsubsection{Hand}\label{d6DcwCJ0jKUAiX5iUpOvw-3}
\\
The \texttt{Hand} class represents a Blackjack hand used in this game and contains multiple cards.

\textbf{The Hand class}

\begin{lstlisting}[language=Java]
import java.util.*;

public class Hand {
    private List<Card> cards;

    public Hand(Card card1, Card card2) { /* ... */ }
    public void addCard(Card card) { /* ... */ }
    public List<Card> getCards() { /* ... */ }
    public int getScore() { /* ... */ }
    @Override
    public String toString() { /* ... */ }
}
\end{lstlisting}

\subsubsection{Players}\label{YlcfyJSbLgicU6bo0jzVP}
\\
The \texttt{Player} is an abstract class, and the \texttt{BlackjackPlayer} and \texttt{Dealer} classes extend the \texttt{Player} class.

\begin{itemize}
\item
\protect\phantomsection\label{2Up2oL0_bCdngdGwQ9Kpl}
\texttt{BlackjackPlayer}: They place the first wager and update the stake with winning and losing sums. They can choose between the hit and stand options.
\item
\protect\phantomsection\label{45wLkzfLSfaxZe8nt114y}
\texttt{Dealer}: They are primarily responsible for dealing cards and following the Blackjack rules.
\end{itemize}

\textbf{Player and its derived classes}

\begin{lstlisting}[language=Java]
public abstract class Player {
    protected String id;
    protected String password;
    protected double balance;
    protected AccountStatus status;
    protected Person person;
    protected Hand hand;

    public Player(String id, String password, double balance, Person person) { /* ... */ }
    public void setHand(Hand hand) { /* ... */ }
    public Hand getHand() { /* ... */ }
    public double getBalance() { /* ... */ }
    public abstract boolean resetPassword();
}

public class BlackjackPlayer extends Player {
    private double bet;

    public BlackjackPlayer(String id, String password, double balance, Person person) { /* ... */ }
    public void placeBet(double amount) { /* ... */ }
    public double getBet() { /* ... */ }
    public void payout(double multiplier) { /* ... */ }
    @Override
    public boolean resetPassword() { /* ... */ }
}

public class Dealer extends Player {
    public Dealer(String id, String password, double balance, Person person) { /* ... */ }
    @Override
    public boolean resetPassword() { /* ... */ }
}
\end{lstlisting}

\subsubsection{Blackjack controller and game view}\label{xuA9-h3Czd3NgpR6lNflw-1}
\\
The \texttt{BlackjackController} class validates the actions (hit or stand) and responds accordingly. The
\texttt{BlackjackGameView} class represents the game view.

\textbf{The BlackjackController and BlackjackGameView classes}

\begin{lstlisting}[language=Java]
public class BlackjackController {
    public boolean validateAction(String action) { /* ... */ }
}

public class BlackjackGameView {
    public void showHands(Hand player, Hand dealer, boolean hideDealerHole) { /* ... */ }
    public void showResult(String result, double payout) { /* ... */ }
}
\end{lstlisting}

\subsubsection{Blackjack game}\label{xuA9-h3Czd3NgpR6lNflw-2}
\\
The \texttt{BlackjackGame} class represents how we can play this game, or its basic sequence of play.

\textbf{The BlackjackGame class}

\begin{lstlisting}[language=Java]
public class BlackjackGame {
    private BlackjackPlayer player;
    private Dealer dealer;
    private Shoe shoe;
    private BlackjackController controller;
    private BlackjackGameView view;

    public BlackjackGame(BlackjackPlayer player, Dealer dealer, int numDecks) { /* ... */ }
    public void start() { /* ... */ }
    public void playAction(String action) { /* ... */ }
    private void hit() { /* ... */ }
    private void stand() { /* ... */ }
    private void compareAndSettle() { /* ... */ }
    private void settle(double multiplier, String message) { /* ... */ }
}
\end{lstlisting}

\subsection{Executable Code: Online Blackjack Game}\label{nkyDasz-Z4uHZJAKq-5m0}
\\
Below is a fully self-contained, runnable program in Java, C\#, C++, Python, and JavaScript that demonstrates the core workflows of the Online Blackjack Game. The main driver code of the system resides in the~\texttt{Driver.java},~\texttt{Driver.cs},~\texttt{Driver.py},~\texttt{Driver.cpp}, and~\texttt{Driver.js}~for all respective languages. You can click the ``Run'' button to execute the code.

\subsubsection{What this code shows}\label{bVgjPIs3owvEWrDwEPetb}

\begin{itemize}
\item
\protect\phantomsection\label{zCD3mlqWNlzHlaeAy_pTn}
\textbf{System initialization:} Sets up the player with an account and balance, a dealer, and a multi-deck shoe.
\item
\protect\phantomsection\label{od_adrgTSjOstOGKT71p2}
\textbf{Scenario 1 (Player gets Blackjack):} Simulates a round where the player achieves Blackjack.
\item
\protect\phantomsection\label{OdXl-SmK1H2dfvI64uO8y}
\textbf{Scenario 2 (Player busts):} Demonstrates a player drawing too many cards and losing their bet.
\item
\protect\phantomsection\label{SSjZvsGZUYPlnAW8_xe-v}
\textbf{Scenario 3 (Dealer busts):} Shows the dealer exceeding 21, allowing the player to win.
\item
\protect\phantomsection\label{WqRM61yLlHltu5Io9Wcrg}
\textbf{Scenario 4 (Push):} Illustrates an outcome where players and dealers have equal scores, resulting in a tie.
\item
\protect\phantomsection\label{PFIcpEIBQ1FN--7Gkha5O}
\textbf{Scenario 5 (Player wins by standing):} Player stands and beats the dealer's hand.
\item
\protect\phantomsection\label{_E5_4VnOlLtDwdncj9lh_}
\textbf{Scenario 6 (Player forfeits):} Player forfeits in the middle of the game and loses their bet.
\end{itemize}
\\
\textbf{Hint: Try customizing the code!}
\\
This code is designed for experimentation, not just reading. You can edit, extend, or modify any part of the example.

\textbf{Executable Code: Online Blackjack Game}

\begin{lstlisting}[language=Java]
import java.util.*;

public class Driver {
    public static void main(String[] args) {
        BlackjackPlayer player = new BlackjackPlayer(1000);
        Dealer dealer = new Dealer();
        BlackjackGame game = new BlackjackGame(player, dealer, 4);

        // Scenario 1: Player gets Blackjack
        System.out.println("
--- Scenario 1: Player gets Blackjack ---");
        System.out.println("
Initial balance $" + player.getBalance());
        player.placeBet(100);
        game.view.announceBet(100, player.getBalance());
        game.start(); // Initial deal
        // Normally, we'd check if player actually got blackjack, but we'll proceed to "stand" to demo flow.
        game.view.announcePlayerAction("stand");
        game.playAction("stand");

        // Scenario 2: Player bust
        System.out.println("
--- Scenario 2: Player bust ---");
        player.placeBet(50);
        game.view.announceBet(50, player.getBalance());
        game.start();
        game.view.announcePlayerAction("hit");
        game.playAction("hit");
        game.view.announcePlayerAction("hit");
        game.playAction("hit"); // Very likely to bust!

        // Scenario 3: Dealer busts
        System.out.println("
--- Scenario 3: Dealer busts ---");
        player.placeBet(75);
        game.view.announceBet(75, player.getBalance());
        game.start();
        game.view.announcePlayerAction("stand");
        game.playAction("stand");

        // Scenario 4: Push (tie)
        System.out.println("
--- Scenario 4: Push (tie) ---");
        player.placeBet(100);
        game.view.announceBet(100, player.getBalance());
        game.start();
        game.view.announcePlayerAction("stand");
        game.playAction("stand");

        // Scenario 5: Player stands & wins
        System.out.println("
--- Scenario 5: Player stands and wins ---");
        player.placeBet(60);
        game.view.announceBet(60, player.getBalance());
        game.start();
        game.view.announcePlayerAction("stand");
        game.playAction("stand");

        // Scenario 6: Player forfeits
        System.out.println("
--- Scenario 6: Forfeit ---");
        player.placeBet(80);
        game.view.announceBet(80, player.getBalance());
        // Forfeit simulation (player leaves):
        game.view.announceForfeit(player.getBalance());
        player.payout(-1); // Loses their bet
    }
}
\end{lstlisting}

\subsection{Wrapping up}\label{oduvT-U5GDbYO8Cu-X2TP}
\\
In this chapter, we've explored the complete design of the Blackjack game. We 've also examined how a basic Blackjack game can be visualized using various UML diagrams and designed using object-oriented principles and design patterns.

% End of section content